% \iffalse meta-comment
% apendix.dtx
% Author: Peter Wilson, Herries Press
% Maintainer: Will Robertson (will dot robertson at latex-project dot org)
% Copyright 1998--2004 Peter R. Wilson
%
% This work may be distributed and/or modified under the
% conditions of the LaTeX Project Public License, either
% version 1.3c of this license or (at your option) any 
% later version: <http://www.latex-project.org/lppl.txt>
%
% This work has the LPPL maintenance status "maintained".
% The Current Maintainer of this work is Will Robertson.
%
% This work consists of the files listed in the README file.
%
% 
%<*driver>
\documentclass{ltxdoc}
\EnableCrossrefs
\CodelineIndex
\setcounter{StandardModuleDepth}{1}
\begin{document}
  \DocInput{appendix.dtx}
\end{document}
%</driver>
%
% \fi
%
% \CheckSum{481}
%
% \DoNotIndex{\',\.,\@M,\@@input,\@addtoreset,\@arabic,\@badmath}
% \DoNotIndex{\@centercr,\@cite}
% \DoNotIndex{\@dotsep,\@empty,\@float,\@gobble,\@gobbletwo,\@ignoretrue}
% \DoNotIndex{\@input,\@ixpt,\@m}
% \DoNotIndex{\@minus,\@mkboth,\@ne,\@nil,\@nomath,\@plus,\@set@topoint}
% \DoNotIndex{\@tempboxa,\@tempcnta,\@tempdima,\@tempdimb}
% \DoNotIndex{\@tempswafalse,\@tempswatrue,\@viipt,\@viiipt,\@vipt}
% \DoNotIndex{\@vpt,\@warning,\@xiipt,\@xipt,\@xivpt,\@xpt,\@xviipt}
% \DoNotIndex{\@xxpt,\@xxvpt,\\,\ ,\addpenalty,\addtolength,\addvspace}
% \DoNotIndex{\advance,\Alph,\alph}
% \DoNotIndex{\arabic,\ast,\begin,\begingroup,\bfseries,\bgroup,\box}
% \DoNotIndex{\bullet}
% \DoNotIndex{\cdot,\cite,\CodelineIndex,\cr,\day,\DeclareOption}
% \DoNotIndex{\def,\DisableCrossrefs,\divide,\DocInput,\documentclass}
% \DoNotIndex{\DoNotIndex,\egroup,\ifdim,\else,\fi,\em,\endtrivlist}
% \DoNotIndex{\EnableCrossrefs,\end,\end@dblfloat,\end@float,\endgroup}
% \DoNotIndex{\endlist,\everycr,\everypar,\ExecuteOptions,\expandafter}
% \DoNotIndex{\fbox}
% \DoNotIndex{\filedate,\filename,\fileversion,\fontsize,\framebox,\gdef}
% \DoNotIndex{\global,\halign,\hangindent,\hbox,\hfil,\hfill,\hrule}
% \DoNotIndex{\hsize,\hskip,\hspace,\hss,\if@tempswa,\ifcase,\or,\fi,\fi}
% \DoNotIndex{\ifhmode,\ifvmode,\ifnum,\iftrue,\ifx,\fi,\fi,\fi,\fi,\fi}
% \DoNotIndex{\input}
% \DoNotIndex{\jobname,\kern,\leavevmode,\let,\leftmark}
% \DoNotIndex{\list,\llap,\long,\m@ne,\m@th,\mark,\markboth,\markright}
% \DoNotIndex{\month,\newcommand,\newcounter,\newenvironment}
% \DoNotIndex{\NeedsTeXFormat,\newdimen}
% \DoNotIndex{\newlength,\newpage,\nobreak,\noindent,\null,\number}
% \DoNotIndex{\numberline,\OldMakeindex,\OnlyDescription,\p@}
% \DoNotIndex{\pagestyle,\par,\paragraph,\paragraphmark,\parfillskip}
% \DoNotIndex{\penalty,\PrintChanges,\PrintIndex,\ProcessOptions}
% \DoNotIndex{\protect,\ProvidesClass,\raggedbottom,\raggedright}
% \DoNotIndex{\refstepcounter,\relax,\renewcommand,\reset@font}
% \DoNotIndex{\rightmargin,\rightmark,\rightskip,\rlap,\rmfamily,\roman}
% \DoNotIndex{\roman,\secdef,\selectfont,\setbox,\setcounter,\setlength}
% \DoNotIndex{\settowidth,\sfcode,\skip,\sloppy,\slshape,\space}
% \DoNotIndex{\symbol,\the,\trivlist,\typeout,\tw@,\undefined,\uppercase}
% \DoNotIndex{\usecounter,\usefont,\usepackage,\vfil,\vfill,\viiipt}
% \DoNotIndex{\viipt,\vipt,\vskip,\vspace}
% \DoNotIndex{\wd,\xiipt,\year,\z@}
%
% \changes{v1.0}{1998/11/29}{First public release}
% \changes{v1.0a}{1999/07/28}{Added text about includes}
% \changes{v1.1}{2000/02/29}{Extended appendices}
% \changes{v1.1}{2000/02/29}{Added subappendices}
% \changes{v1.1a}{2001/03/15}{Fixed problem with \cs{addappheadtotoc}}
% \changes{v1.2}{2002/08/06}{Compatibility with hyperref bookmarks}
% \changes{v1.2}{2002/08/06}{Don't need the ifthen package any more}
% \changes{v1.2a}{2004/04/16}{Changed license and contact details}
% \changes{v1.2b}{2009/09/02}{New maintainer (Will Robertson)}
%
% \def\dtxfile{appendix.dtx}
% \def\fileversion{v1.1}  \def\filedate{2000/02/29}
% \def\fileversion{v1.1a} \def\filedate{2001/03/15}
% \def\fileversion{v1.2}  \def\filedate{2002/08/06}
% \def\fileversion{v1.2a} \def\filedate{2004/04/16}
% \def\fileversion{v1.2b} \def\filedate{2009/09/02}
%
% \newcommand*{\Lpack}[1]{\textsf {#1}}           ^^A typeset a package
% \newcommand*{\Lopt}[1]{\textsf {#1}}            ^^A typeset an option
% \newcommand*{\file}[1]{\texttt {#1}}            ^^A typeset a file
% \newcommand*{\Lcount}[1]{\textsl {\small#1}}    ^^A typeset a counter
% \newcommand*{\pstyle}[1]{\textsl {#1}}          ^^A typeset a pagestyle
% \newcommand*{\Lenv}[1]{\texttt {#1}}            ^^A typeset an environment
%
% \title{The \Lpack{appendix} package\thanks{This
%        file (\texttt{\dtxfile}) has version number \fileversion, last revised
%        \filedate.}}
%
% \author{
%   Author: Peter Wilson, Herries Press\\
%   Maintainer: Will Robertson\\
%   \texttt{will dot robertson at latex-project dot org}
% }
% \date{\filedate}
% \maketitle
% \begin{abstract}
%    The \Lpack{appendix} package provides some facilities for modifying the
% typesetting of appendix titles. Further, |(sub)appendices| environments
% are available that can be used, for example, for per chapter/section
% appendices.
%
% The package
% is designed to work only with classes that have a |\chapter| and/or
% |\section| command. It has not been tested with other packages
% that change the definitions of the sectioning commands. 
% \end{abstract}
% \tableofcontents
%
%
%
%
% \section{Introduction}
%
% In the standard classes the |\appendix| command does the following:
% \begin{itemize}
% \item For classes with Chapters:
%    \begin{itemize}
%    \item Resets the chapter and section counters to zero
%    \item Sets |\@chapapp| to |\appendixname|.
%    \item Redefines |\thechapter| to produce alphabetic appendix numbers.
%    \end{itemize}
% \item For the other classes
%    \begin{itemize}
%    \item Resets the section and subsection counters to zero.
%    \item Redefines |\thesection| to produce alphabetic appendix numbers.
%    \end{itemize}
%  \end{itemize}
%
% The \Lpack{appendix} package provides additional appendixing capabilities.
% It cooperates with the \Lpack{hyperref} 
% package\footnote{With thanks to Hylke W. van Dijk 
% (\texttt{hylke@ubicom.tudelft.nl}) who pointed out that version~1.1 did 
% not and set me on the track for supporting the \Lpack{hyperref} package.}
% but may be problematic when used with packages that change the definition
% of the sectioning commands.
%
% Portions of the package were developed as part of a class
% and package bundle for typesetting ISO standards~\cite{PRW96i}.
% This manual is typeset according to the conventions of the
% \LaTeX{} \textsc{docstrip} utility which enables the automatic
% extraction of the \LaTeX{} macro source files~\cite{GOOSSENS94}.
%
%    Section~\ref{sec:usc} describes the usage of the package.
% Commented source code for the package is in Section~\ref{sec:code}.
%
% \section{The \Lpack{appendix} package} \label{sec:usc}
%
% The \Lpack{appendix} package provides some commands that can be
% used in addition to the |\appendix| command. It also provides
% a new environment that can be used instead of the |\appendix| command.
% The environment provides some addtional actions with respect to
% the simple |\appendix|. First the new commands will be described and
% then the new environment will be discussed.
%
% \DescribeMacro{\appendixpage}
%  The |\appendixpage| command will typeset a heading in the style of
% a |\part| heading for the class. The wording of the heading is
% given by the value of |\appendixpagename|.
%
% \DescribeMacro{\addappheadtotoc}
% This command will insert general heading into the Table of Contents (ToC).
% The text is given by the value of |\appendixtocname|. If used, the command
% must come before the first appendix, as it is meant to be used to introduce
% the appendix titles in the ToC.
%
% \changes{v1.1a}{2001/03/15}{Added NOTE about \cs{addappheadtotoc} page number}
% The above commands can be used in conjunction with the traditional
% |\appendix| command, which they should immediately follow. For example:
% \begin{verbatim}
% \appendix
% \appendixpage
% \addappheadtotoc
% \end{verbatim}
%
% \DescribeMacro{\noappendicestocpagenum}
% \DescribeMacro{\appendicestocpagenum}
% By default the |\addappheadtotoc| command puts a page number
% in the ToC. This can be prevented by using |\noappendicestocpagenum|.
% For symmetry, the |\appendicestocpagenum| command ensures that
% a page number is put in the ToC.
%
% \textbf{NOTE:} Unless |\noappendicestocpagenum| is used 
% the |\addappheadtotoc| command uses the 
% current page number
% when it makes the entry in the ToC. The |\appendixpage| command puts
% a heading in the document like a |\part| heading; in un-chaptered documents
% the |\part| heading appears in the ordinary run of the text like a 
% |\section| heading, but in chaptered documents it is on a page by itself.
% That is, in chaptered documents |\appendixpage| does a |\clear[double]page|
% typesets the heading, and then does another |\clear[double]page|. Therefore,
% in a chaptered document the above sequence of 
% commands will use the page number
% \emph{after} the |\appendixpage| as the ToC 
% entry\footnote{With thanks to Eduardo Jacob (\texttt{edu@kender.es})
% for pointing this out.} and if the ordering
% is reversed (i.e., |\addappheadtotoc| |\appendixname|) then the page number
% \emph{before} |\appendixname| will be used as the ToC entry. For chaptered
% documents it is probably best to do: \\
% |    \clearpage % or \cleardoublepage| \\
% |    \addappheadtotoc| \\
% |    \appendixpage| \\
% which will use the page number of |\appendixpage| as the ToC entry.
% 
%
% \DescribeMacro{\appendixname}
% \DescribeMacro{\appendixtocname}
% \DescribeMacro{\appendixpagename}
% The |\appendixname| command is defined in those classes that provide
% for chapters. It is provided in this package whether or not it has
% been defined in the class. It's default value is `Appendix'.
% The default value of both |\appendixtocname| and |\appendixpagename| is 
% `Appendices'. These names can all be changed via |\renewcommand|.
% For example,
% \begin{verbatim}
% \renewcommand{\appendixtocname}{List of appendices}
% \end{verbatim}
%
% \DescribeEnv{appendices}
% The |appendices| environment can be used instead of the |\appendix|
% command. It provides more functionality than is possible from the 
% combination of the |\appendix|, |\addappheadtotoc| and |\appendixpage|
% commands.
% The functions of the |appendices| environment are usually accessed through
% the package options, but there are declarations that mey be used insted. 
% The options are:
% \begin{itemize}
% \item \Lopt{toc} Put a header (e.g., `Appendices') into the Table
%       of Contents (the ToC) before listing the appendices. (This
%       is done by calling the |\addappheadtotoc| command.)
% \item \Lopt{page} Puts a title (e.g., `Appendices') into the document
%       at the point where the |appendices| environment is begun.
%       (This is done by calling the |\appendixpage| command.)
% \item \Lopt{title} Adds a name (e.g., `Appendix') before each appendix
%       title in the body of the document. 
%       The name is given by the value of |\appendixname|.
%       Note that this is the default behaviour for classes that have
%       chapters.
% \item \Lopt{titletoc} Adds a name (e.g., `Appendix') before each
%        appendix listed in the ToC.
%       The name is given by the value of |\appendixname|.
% \item \Lopt{header} Adds a name (e.g., `Appendix') before each appendix
%        in page headers.
%       The name is given by the value of |\appendixname|.
%       Note that this is the default behaviour for classes that have
%       chapters.
% \end{itemize}
%
%    Depending on the particular package options that are set and the
% document class, the
% |appendices| environment may change the definition of elements of
% the sectioning commands (e.g., |\chapter| or |\section|).
% This may be a problem if the environment
% is used in conjunction with any other package that makes changes to
% these commands. If this is the case, then you will have to examine the
% code for the |appendices| environment and make any necessary changes
% to one or the other of the packages (via your own package file).
% The changes to the sectional heading commands are discarded at the
% end of the |appendices| environment.
%
% \DescribeMacro{\appendixtocon}
% \DescribeMacro{\appendixtocoff}
% |\appendixtocon| is a declaration equivalent to the \Lopt{toc} option.
% The |\appendixtocoff| declaration is equivalent to not using that option.
%
% \DescribeMacro{\appendixpageon}
% \DescribeMacro{\appendixpageoff}
% |\appendixpagecon| is a declaration equivalent to the \Lopt{page} option.
% The |\appendixpageoff| declaration is equivalent to not using 
% that option.
%
% \DescribeMacro{\appendixtitleon}
% \DescribeMacro{\appendixtitleoff}
% |\appendixtitleon| is a declaration equivalent to the \Lopt{title} option.
% The |\appendixtitleoff| declaration is equivalent to not using 
% that option.
%
% \DescribeMacro{\appendixtitletocon}
% \DescribeMacro{\appendixtitletocoff}
% |\appendixtitletocon| is a declaration equivalent to the \Lopt{titletoc} option.
% The |\appendixtitletocoff| declaration is equivalent to not using 
% that option.
%
% \DescribeMacro{\appendixheaderon}
% \DescribeMacro{\appendixheaderoff}
% |\appendixheaderon| is a declaration equivalent to the \Lopt{header} option.
% The |\appendixheaderoff| declaration is equivalent to not using 
% that option.
%
% \DescribeMacro{\restoreapp}
%     The |appendices| environment restores the prior value of the 
% chapter/section counter at the end of the environment, so the
% environment may be used between the main document divisions. By default,
% the appendix counter value is saved and restored by the environment. That
% means that appendices in a series of |appendices| environments will
% be lettered sequentially. To make the lettering start from A each time,
% put the following into the preamble: \\
% |\renewcommand{\restoreapp}{}|
%
% \DescribeEnv{subappendices}
% Within the |subappendices| environment, an appendix is introduced by a
% |\section| command in chaptered documents, otherwise it is introduced
% by a |\subsection| command. Effectively, this provides for appendices 
% at the end of a main document division, as an integral
% part of the division. The |subappendices| environment supports only
% the \Lopt{title} and \Lopt{titletoc} options.
%
% \DescribeMacro{\setthesection}
% \DescribeMacro{\setthesubsection}
% By default, the `subappendices' are numbered like normal (sub)sections,
% except that the (sub)section number itself is typeset as an uppercase
% letter. This behaviour can be changed by redefining these |\setthe...|
% commands. For example, to just have a letter not prepended by the main
% division number, do: \\
% |\renewcommand{\setthesection}{\Alph{section}}| or \\
% |\renewcommand{\setthesubsection}{\Alph{subsection}}| as appropriate.
%
% \subsection{Known problems}
%
%    There is an unfortunate interaction between the \LaTeX{} kernel commands
% |\include| and |\addcontentsline|. If these are used like this:
% \begin{verbatim}
% \addcontentsline{toc}{...}{addtotoc}
% \documentclass[DIV=8, parskip=half, pagesize=auto]{scrartcl}

\usepackage{fixltx2e}
\usepackage{etex}
\usepackage{xspace}
\usepackage{lmodern}
\usepackage[T1]{fontenc}
\usepackage{textcomp}
\usepackage{booktabs}
\usepackage{microtype}
\usepackage{hyperref}

\newcommand*{\mail}[1]{\href{mailto:#1}{\texttt{#1}}}
\newcommand*{\pkg}[1]{\textsf{#1}}
\newcommand*{\cls}[1]{\textsf{#1}}
\newcommand*{\cs}[1]{\texttt{\textbackslash#1}}
\makeatletter
\newcommand*{\cmd}[1]{\cs{\expandafter\@gobble\string#1}}
\makeatother
\newcommand*{\env}[1]{\texttt{#1}}
\newcommand*{\opt}[1]{\texttt{#1}}
\newcommand*{\meta}[1]{\textlangle\textsl{#1}\textrangle}
\newcommand*{\marg}[1]{\texttt{\{}\meta{#1}\texttt{\}}}
\newcommand*{\oarg}[1]{\texttt{[}\meta{#1}\texttt{]}}
\newcommand*{\BibTeX}{Bib\kern-0.08em\TeX\@\xspace}
\newcommand*{\BigTeX}{Big\kern-0.08em\TeX\@\xspace}

\pdfstringdefDisableCommands{%
  \def\BibTeX{BibTeX\xspace}%
  \def\BigTeX{BigTeX\space}%
}

\addtokomafont{title}{\rmfamily}

\title{The \pkg{import} package\thanks{This manual corresponds to \pkg{import}~v5.1, dated~23--Mar--2009.}}
\author{Donald Arseneau  (\mail{asnd@triumf.ca})}
\date{23--Mar--2009}


\begin{document}

\maketitle

This software is in the public domain; free of any restrictions.

Two new \LaTeX\ commands, ``\cmd{\import}\marg{full\_path}\marg{file}''\\ and
``\cmd{\subimport}\-\marg{path\_extension}\marg{file}'' are defined to input a file 
from another directory, allowing that file to find its own inputs
(using ``\cmd{\input}'', ``\cmd{\includegraphics}'' etc.)\ in that directory.

Alias command names are ``\cmd{\inputfrom}'' and ``\cmd{\subinputfrom}''.

Also provided are ``\cmd{\includefrom}'' and ``\cmd{\subincludefrom}'', which are
based on the ``\cmd{\include}'' command, rather than ``\cmd{\input}''.  There are
also ``\texttt{*}'' variants described below.

For example, if a remote file ``\texttt{/usr5/friend/work/report.tex}'' has contents:
%
\begin{verbatim}
My graph: \includegraphics{picture}
\input{explanation}
\end{verbatim}
%
then you can safely input that file in your own document with the command
``\verb+\import{/usr5/friend/work/}{report}+''; the explanation and picture will 
be taken from the ``\texttt{/usr5/friend/work/}'' directory.

The ``\cmd{\subimport}'' command takes a relative path instead of a full absolute
file path, and it allows imported files to import files themselves, using
their own directory as the root of another ``\meta{path\_extension}''. Do not use
both ``\cmd{\import}'' and ``\cmd{\subimport}'' in the same file.

For example, if a file is imported (using either command) from directory
``\texttt{abc/}'' and that file contains the command ``\verb+\subimport{lmn/}{xyz}+'' then
file ``\texttt{abc/lmn/xyz.tex}'' is input, and any ``\cmd{\input}'' commands in that file
will read files from directory ``\texttt{abc/lmn/}''. 

Note that the sub-import path is merely appended to the current import
path.  Mistakes from this method must be rectified by ``\cmd{\import@path@fix}''.

Depending on on how your \TeX\ system is configured, if a file with the 
same name as the import file exists in the current directory or in the 
\texttt{TEXINPUTS} path, that other file will be read in preference to one in the 
import directory.  So here is the real behavior of the previous example:
Given the sequence ``\verb+\import{abc/}{one}+'', ``\verb+\subimport{lmn/uvw/}{two}+'' (in
file \texttt{one}), ``\verb+\input{three}+'' (in file \texttt{two}), \LaTeX\ first looks for \texttt{three.tex}
in the \texttt{TEXINPUTS} search path; if not found, it tries ``\texttt{abc/lmn/uvw/three}'';
if that doesn't exist, it tries ``\texttt{abc/three}''; if still not found, it tries
the defined ``\cmd{\input@path}'', if there is one.

To avoid searching the \texttt{TEXINPUTS} path when importing files, use the `star'
versions of the commands: ``\cmd{\import*}'' and ``\cmd{\subimport*}''.

A hook ``\cmd{\import@path@fix}'' is provided to reformat the import path
to fit the syntax of a particular operating system.  It \emph{could} be
defined to convert unix-style path names to the local format, but 
all it does now is remove ``\texttt{][}'' from VMS sub-import directories.

Presently, the paths are defined ``locally'' so input files must have 
balanced grouping.

\end{document}

% \end{verbatim}
% then the text of the |\addcontentsline| command (`addtotoc' in the example)
% is not written to the
% appropriate (toc) file until \textit{after} the included file has written all
% its entries out to the (toc) file. As far as I can tell, there is no way 
% around this behaviour without rewriting parts of the \LaTeX{} kernel code.
% 
%     It is thus up to the author to avoid putting an |\addcontentsline| command
% (or a command that internally uses |\addcontentsline|, as does the
% |\addappheadtotoc| command) before
% an |\include|d file that writes out to the same file. Things work as
% expected if the |\addcontentsline| command is placed within the |\include|d
% file, or if the imported file is |\input|ed instead of |\include|d.
%
% \StopEventually{}
%
% \section{The package code} \label{sec:code}
%
%    Announce the name and version of the package, which requires
% \LaTeXe.
%    \begin{macrocode}
%<*usc>
\NeedsTeXFormat{LaTeX2e}
\ProvidesPackage{appendix}[2009/09/02 v1.2b extra appendix facilities]

%    \end{macrocode}
%
% In order to try and avoid name clashes with other packages, each internal
% name will include the character string |@pp|.
%
%
% \begin{macro}{\if@knownclass@pp}
% \begin{macro}{\if@chapter@pp}
%    These are used when we need to decide what appendix style is being used
%    for the document. Assume the \Lpack{article} class or other without
% chapters.
% \changes{v1.1a}{2001/03/15}{Checking on sectional commands, not classes}
%    \begin{macrocode}
\newif\if@chapter@pp\@chapter@ppfalse
\newif\if@knownclass@pp\@knownclass@ppfalse
%    \end{macrocode}
%    Check the sectioning commands.
%    \begin{macrocode}
\@ifundefined{chapter}{%
  \@ifundefined{section}{}{\@knownclass@pptrue}}{%
  \@chapter@pptrue\@knownclass@pptrue}
%    \end{macrocode}
% \end{macro}
% \end{macro}
%
% \begin{macro}{\phantomsection}
% \begin{macro}{\the@pps}
% \begin{macro}{\if@pphyper}
% We need to provide |\phantomsection| if \Lpack{hyperref} is not
% used and, whether or not \Lpack{hyperref} is used, we need to define
% a counter here to support potential hyperrefs (used to disambiguate
% (sub)appendices).
% |\if@pphyper| is TRUE if the \Lpack{hyperref} package is used.
% \changes{v1.2}{2002/08/06}{Added the \texttt{@pps} counter}
% \changes{v1.2}{2002/08/06}{Added \cs{if@pphyper}}
%    \begin{macrocode}
\providecommand{\phantomsection}{}
\newcounter{@pps}
  \renewcommand{\the@pps}{\alph{@pps}}
\newif\if@pphyper
  \@pphyperfalse
\AtBeginDocument{%
  \@ifpackageloaded{hyperref}{\@pphypertrue}{}}

%    \end{macrocode}
% \end{macro}
% \end{macro}
% \end{macro}
%
% \begin{macro}{\if@dotoc@pp}
% \begin{macro}{\if@dotitle@pp}
% \begin{macro}{\if@dotitletoc@pp}
% \begin{macro}{\if@dohead@pp}
% \begin{macro}{\if@dopage@pp}
%    A set of booleans for the options. Default is the |appendices|
% environment does nothing more than the |\appendix| command does
% unless one or more options are set.
%    \begin{macrocode}
\newif\if@dotoc@pp\@dotoc@ppfalse
\newif\if@dotitle@pp\@dotitle@ppfalse
\newif\if@dotitletoc@pp\@dotitletoc@ppfalse
\newif\if@dohead@pp\@dohead@ppfalse
\newif\if@dopage@pp\@dopage@ppfalse
%    \end{macrocode}
% \end{macro}
% \end{macro}
% \end{macro}
% \end{macro}
% \end{macro}
%
%    Now we can do the five options.
%    \begin{macrocode}
\DeclareOption{toc}{\@dotoc@pptrue}
\DeclareOption{title}{\@dotitle@pptrue}
\DeclareOption{titletoc}{\@dotitletoc@pptrue}
\DeclareOption{header}{\@dohead@pptrue}
\DeclareOption{page}{\@dopage@pptrue}
%    \end{macrocode}
% Process the options now.
%    \begin{macrocode}
\ProcessOptions\relax
%    \end{macrocode}
%
% Issue a warning if |\chapter| and |\section| are undefined, then
% quit.
%    \begin{macrocode}
\newcommand{\@ppendinput}{}
\if@knownclass@pp\else
  \PackageWarningNoLine{appendix}%
    {There is no \protect\chapter\space or \protect\section\space command.\MessageBreak
     The appendix package will not be used}
  \renewcommand{\@ppendinput}{\endinput}
\fi
\@ppendinput

%    \end{macrocode}
%
% \begin{macro}{\appendixtocon}
% \begin{macro}{\appendixtocoff}
% Declarative forms of the \Lopt{toc} option.
% \changes{v1.2}{2002/08/06}{Added declarations for the options}
%    \begin{macrocode}
\newcommand{\appendixtocon}{\@dotoc@pptrue}
\newcommand{\appendixtocoff}{\@dotoc@ppfalse}
%    \end{macrocode}
% \end{macro}
% \end{macro}
%
% \begin{macro}{\appendixpageon}
% \begin{macro}{\appendixpageoff}
% Declarative forms of the \Lopt{page} option.
%    \begin{macrocode}
\newcommand{\appendixpageon}{\@dopage@pptrue}
\newcommand{\appendixpageoff}{\@dopage@ppfalse}
%    \end{macrocode}
% \end{macro}
% \end{macro}
%
% \begin{macro}{\appendixtitleon}
% \begin{macro}{\appendixtitleoff}
% Declarative forms of the \Lopt{title} option.
%    \begin{macrocode}
\newcommand{\appendixtitleon}{\@dotitle@pptrue}
\newcommand{\appendixtitleoff}{\@dotitle@ppfalse}
%    \end{macrocode}
% \end{macro}
% \end{macro}
%
% \begin{macro}{\appendixtitletocon}
% \begin{macro}{\appendixtitletocoff}
% Declarative forms of the \Lopt{titletoc} option.
%    \begin{macrocode}
\newcommand{\appendixtitletocon}{\@dotitletoc@pptrue}
\newcommand{\appendixtitletocoff}{\@dotitletoc@ppfalse}
%    \end{macrocode}
% \end{macro}
% \end{macro}
%
% \begin{macro}{\appendixheaderon}
% \begin{macro}{\appendixheaderoff}
% Declarative forms of the \Lopt{header} option.
%    \begin{macrocode}
\newcommand{\appendixheaderon}{\@dohead@pptrue}
\newcommand{\appendixheaderoff}{\@dohead@ppfalse}
%    \end{macrocode}
% \end{macro}
% \end{macro}
%
% \begin{macro}{\@ppsavesec}
% \begin{macro}{\@pprestoresec}
% \begin{macro}{\@ppsaveapp}
% \begin{macro}{\restoreapp}
%  For the |appendices| environment we need to save and restore the
% main document division number and the appendix number. The |\restoreapp|
% command is the one for the user.
% \changes{v1.1}{2000/02/29}{Added commands to save and restore sectional numbering}
%    \begin{macrocode}
\newcounter{@ppsavesec}
\newcounter{@ppsaveapp}
\setcounter{@ppsaveapp}{0}
\newcommand{\@ppsavesec}{%
  \if@chapter@pp \setcounter{@ppsavesec}{\value{chapter}} \else
                 \setcounter{@ppsavesec}{\value{section}} \fi}
\newcommand{\@pprestoresec}{%
  \if@chapter@pp \setcounter{chapter}{\value{@ppsavesec}} \else
                 \setcounter{section}{\value{@ppsavesec}} \fi}
\newcommand{\@ppsaveapp}{%
  \if@chapter@pp \setcounter{@ppsaveapp}{\value{chapter}} \else
                 \setcounter{@ppsaveapp}{\value{section}} \fi}
\newcommand{\restoreapp}{%
  \if@chapter@pp \setcounter{chapter}{\value{@ppsaveapp}} \else
                 \setcounter{section}{\value{@ppsaveapp}} \fi}
%    \end{macrocode}
% \end{macro}
% \end{macro}
% \end{macro}
% \end{macro}
%
% \begin{macro}{\appendixname}
% \begin{macro}{\appendixtocname}
% \begin{macro}{\appendixpagename}
%  These commands hold the names that might be used. |\appendixname|
% may have been defined in the class. The others are new.
%    \begin{macrocode}
\providecommand{\appendixname}{Appendix}
\newcommand{\appendixtocname}{Appendices}
\newcommand{\appendixpagename}{Appendices}
%    \end{macrocode}
% \end{macro}
% \end{macro}
% \end{macro}
%
% \begin{macro}{\appendixpage}
% The command to typeset a page announcing the start of the appendices.
% It is based on the |\part| definition (either from the \Lpack{book}
% class or the \Lpack{article} class). 
%    \begin{macrocode}
\newcommand{\appendixpage}{%
  \if@chapter@pp \@chap@pppage \else \@sec@pppage \fi
}
%    \end{macrocode}
% \end{macro}
%
% \begin{macro}{\clear@ppage}
%  The non-chaptered classes do not define |\if@openright|, but we need to 
% use this for chaptered documents to clear the appropriate pages.
% |\clear@ppage| does the right thing, but must only be called in
% chapter related code, otherwise there will be error message like 
% |extra \else| or |extra \fi|.
%    \begin{macrocode}
\newcommand{\clear@ppage}{%
  \if@openright\cleardoublepage\else\clearpage\fi}

%    \end{macrocode}
% \end{macro}
%
% \begin{macro}{\@chap@pppage}
% Do an appendix page in chapter style.
% Copy code from the \Lpack{book} class |\part| command, but use 
% |\appendixpagename| as the title.
%    \begin{macrocode}
\newcommand{\@chap@pppage}{%
  \clear@ppage
  \thispagestyle{plain}%
  \if@twocolumn\onecolumn\@tempswatrue\else\@tempswafalse\fi
  \null\vfil
  \markboth{}{}%
  {\centering
   \interlinepenalty \@M
   \normalfont
   \Huge \bfseries \appendixpagename\par}%
%    \end{macrocode}
% Add to ToC if requested
%    \begin{macrocode}
  \if@dotoc@pp
    \addappheadtotoc
  \fi
%    \end{macrocode}
% In the \Lpack{book} class the |\part| command is finished off by calling
% |\@endpart|. There are two problems with this in this package. (1)
% |\@endpart| is not defined in \Lpack{article} style classes and (2)
% it always throws a blank page which does not look good if the \Lopt{openany}
% option is used. So, code it all up here.
%    \begin{macrocode}
  \vfil\newpage
  \if@twoside
    \if@openright
      \null
      \thispagestyle{empty}%
      \newpage
    \fi
  \fi
  \if@tempswa
    \twocolumn
  \fi
}

%    \end{macrocode}
% \end{macro}
%
% \begin{macro}{\@sec@pppage}
% Copy code from the \Lpack{article} class |\part| command, but use 
% |\appendixpagename| as 
% the title.
%    \begin{macrocode}
\newcommand{\@sec@pppage}{%
  \par
  \addvspace{4ex}%
  \@afterindentfalse
  {\parindent \z@ \raggedright
   \interlinepenalty \@M
   \normalfont
   \huge \bfseries \appendixpagename%
   \markboth{}{}\par}%
%    \end{macrocode}
% Add to ToC if requested
%    \begin{macrocode}
  \if@dotoc@pp
    \addappheadtotoc
  \fi
  \nobreak
  \vskip 3ex
  \@afterheading
}

%    \end{macrocode}
% \end{macro}
%
% \begin{macro}{\if@pptocpage}
% \begin{macro}{\noappendicestocpagenum}
% \begin{macro}{\appendicestocpagenum}
% \begin{macro}{\addappheadtotoc}
% The |\addappheadtotoc| command adds an `appendices' line to the ToC. 
% The style is the same
% as used in \Lpack{tocbibind} for the `List of figures' line. That is,
% as a Chapter heading or a Section heading. |\if@pptocpage| controls
% whether ot not a page number is put into the ToC.
% \changes{v1.2}{2002/08/06}{Added \cs{noappendicestocpagenum} and changed
%                           \cs{addappheadtotoc}}
%    \begin{macrocode}
\newif\if@pptocpage
  \@pptocpagetrue
\newcommand{\noappendicestocpagenum}{\@pptocpagefalse}
\newcommand{\appendicestocpagenum}{\@pptocpagetrue}
\newcommand{\addappheadtotoc}{%
  \phantomsection
  \if@chapter@pp
%    \end{macrocode}
% Chaptered document
%    \begin{macrocode}
    \if@pptocpage
      \addcontentsline{toc}{chapter}{\appendixtocname}%
    \else
      \if@pphyper
        \addtocontents{toc}%
          {\protect\contentsline{chapter}{\appendixtocname}{}{\@currentHref}}%
      \else
        \addtocontents{toc}%
          {\protect\contentsline{chapter}{\appendixtocname}{}}%
      \fi
    \fi      
  \else
%    \end{macrocode}
% Not a chaptered document
%    \begin{macrocode}
    \if@pptocpage
      \addcontentsline{toc}{section}{\appendixtocname}%
    \else
      \if@pphyper
        \addtocontents{toc}%
          {\protect\contentsline{section}{\appendixtocname}{}{\@currentHref}}%
      \else
        \addtocontents{toc}%
          {\protect\contentsline{section}{\appendixtocname}{}}%
      \fi
    \fi
  \fi
}

%    \end{macrocode}
% \end{macro}
% \end{macro}
% \end{macro}
% \end{macro}
%
% For my reference,
% here is the standard version of the |\appendix| macro, but modified for
% both chaptered and unchaptered documents.
% \begin{verbatim}
% \newcommand{\appendix}{\par
%   \if@chapter@pp
%     \setcounter{chapter}{0}%
%     \setcounter{section}{0}%
%     \gdef\@chapapp{\appendixname}%
%     \gdef\thechapter{\@Alph\c@chapter}
%   \else
%     \setcounter{section}{0}%
%     \setcounter{subsection}{0}%
%     \gdef\thesection{\@Alph\c@section}
%   \fi
% }
% \end{verbatim}
%
% And this equivalently is what the \Lpack{hyperref} package does.
% \begin{verbatim}
% \def\Hy@chapterstring{chapter}
% \def\Hy@appendixstring{appendix}
% \def\Hy@chapapp{\Hy@chapterstring}
% \let\Hy@org@appendix\appendix
% \def\appendix{%
%    \Hy@org@appendix
%    \if@chapter@pp
%      \gdef\theHchapter{\Alph{chapter}}%
%    \else
%      \gdef\theHsection{\Alph{section}}%
%    \fi
%    \xdef\Hy@chapapp{\Hy@appendixstring}%
% }
% \end{verbatim}
%
% \begin{macro}{\theH@pps}
% We are going to use |\theH@pps| to disambiguate contents of appendices
% that might have the same hyperref marks. It is |\provide|d as if 
% the \Lpack{appendix} and \Lpack{hyperref} are in the `wrong' order
% then somehow \Lpack{hyperref} defines it before \Lpack{appendix}
% can get to it.
% \changes{v1.2}{2002/08/06}{Added \cs{theH@pps}}
%    \begin{macrocode}
\providecommand{\theH@pps}{\alph{@pps}}

%    \end{macrocode}
% \end{macro}
%
% \begin{macro}{\@resets@pp}
% Resets the appropriate sectioning counters and names. This does almost
% exactly
% what the default |\appendix| command does, except that it saves and 
% restores sectional numbering. It saves the sectional number at the start
% and restores the appendix number at the end.
% \changes{v1.1}{2000/02/29}{Added number save/restore to \cs{@reset@pp}}
% \changes{v1.2}{2002/08/06}{Added hyperref code to \cs{@reset@pp}}
%    \begin{macrocode}
\newcommand{\@resets@pp}{\par
  \@ppsavesec
  \stepcounter{@pps}
  \setcounter{section}{0}%
  \if@chapter@pp
    \setcounter{chapter}{0}%
    \renewcommand\@chapapp{\appendixname}%
    \renewcommand\thechapter{\@Alph\c@chapter}%
  \else
    \setcounter{subsection}{0}%
    \renewcommand\thesection{\@Alph\c@section}%
  \fi
  \if@pphyper
%    \end{macrocode}
% Now handle the \Lpack{hyperref} tweaks.
%    \begin{macrocode}
    \if@chapter@pp
      \renewcommand{\theHchapter}{\theH@pps.\Alph{chapter}}%
    \else
      \renewcommand{\theHsection}{\theH@pps.\Alph{section}}%
    \fi
    \def\Hy@chapapp{\appendixname}%
  \fi
  \restoreapp
}

%    \end{macrocode}
% \end{macro}
%
% \begin{environment}{appendices}
%  This is the heart of the package. Start it off by doing the resetting
% done by the |\appendix| command. Then do the simple options before
% getting into the complications of redefinitions. Remember to take care
% of an interaction between |\addappheadtotoc| and |\appendixpage|.
% \changes{v1.1a}{2001/03/15}{Changed implementation of easy options in \texttt{appendices} environment}
%    \begin{macrocode}
\newenvironment{appendices}{%
  \@resets@pp
  \if@dotoc@pp 
    \if@dopage@pp              % both page and toc
      \if@chapter@pp           % chapters
        \clear@ppage
      \fi
      \appendixpage
    \else                      % toc only
       \if@chapter@pp          % chapters
         \clear@ppage
       \fi
      \addappheadtotoc
    \fi
  \else
    \if@dopage@pp              % page only
      \appendixpage
    \fi
  \fi
%    \end{macrocode}
% There is only one other option applicable to the chapter style, so do
% it now and clear the way for doing the section style. To implement
% the \Lopt{titletoc} option, we redefine the |\addcontentsline| command.
%    \begin{macrocode}
  \if@chapter@pp
    \if@dotitletoc@pp \@redotocentry@pp{chapter} \fi
  \else
%    \end{macrocode}
% The rest of the code is specific to the section style. While we're in the 
% mood we might as well finish off doing the \Lopt{titletoc} option.
%    \begin{macrocode}
    \if@dotitletoc@pp \@redotocentry@pp{section} \fi
%    \end{macrocode}
% The next piece of code implements the \Lopt{header} option by providing
% a special version of |\sectionmark|.
%    \begin{macrocode}
    \if@dohead@pp 
      \def\sectionmark##1{%
        \if@twoside
          \markboth{\@formatsecmark@pp{##1}}{}
        \else
          \markright{\@formatsecmark@pp{##1}}{}
        \fi}
    \fi
%    \end{macrocode}
% The next piece of code implements the \Lopt{title} option by doing cunning
% things with the |\@seccntformat|.\footnote{From a posting to 
% \texttt{comp.tex.tex} by Donald Arseneau on 13 August 1998.}
%    \begin{macrocode}
    \if@dotitle@pp
      \def\sectionname{\appendixname}
      \def\@seccntformat##1{\@ifundefined{##1name}{}{\csname ##1name\endcsname\ }%
        \csname the##1\endcsname\quad}
    \fi
  \fi}{%
%    \end{macrocode}
% At the end of the environment, save the appendix number and restore the
% sectional number.
% \changes{v1.1}{2000/02/29}{Changed end of appendix environment}
%    \begin{macrocode}
  \@ppsaveapp\@pprestoresec}

%    \end{macrocode}
% \end{environment}
%
% \begin{macro}{\setthesection}
% \begin{macro}{\setthesubsection}
% The user commands for specifying the numbering style for subappendices.
% \changes{v1.1}{2000/02/29}{Added \cs{setthesection} and \cs{setthesubsection} commands}
%    \begin{macrocode}
\newcommand{\setthesection}{\thechapter.\Alph{section}}
\newcommand{\setthesubsection}{\thesection.\Alph{subsection}}

%    \end{macrocode}
% \end{macro}
% \end{macro}
%
% \begin{macro}{\@resets@ppsub}
% Similar to |\@resets@pp| except that it is for use within the 
% |subappendices| envirionment; as such, it is a bit simpler.
% \changes{v1.1}{2000/02/29}{Added \cs{@resets@ppsub} command}
% \changes{v1.2}{2002/08/07}{Added hyperref code to \cs{@resets@ppsub}}
%    \begin{macrocode}
\newcommand{\@resets@ppsub}{\par
  \stepcounter{@pps}
  \if@chapter@pp
    \setcounter{section}{0}
    \renewcommand{\thesection}{\setthesection}
  \else
    \setcounter{subsection}{0}
    \renewcommand{\thesubsection}{\setthesubsection}
  \fi
  \if@pphyper
%    \end{macrocode}
% Now handle the \Lpack{hyperref} tweaks.
%    \begin{macrocode}
    \if@chapter@pp
      \renewcommand{\theHsection}{\theH@pps.\setthesection}%
    \else
      \renewcommand{\theHsubsection}{\theH@pps.\setthesubsection}%
    \fi
    \def\Hy@chapapp{\appendixname}%
  \fi
}

%    \end{macrocode}
% \end{macro}
%
% \begin{environment}{subappendices}
%  The environment for subappendices. Start it off by doing the resetting
% of the |\(sub)section| command. 
% \changes{v1.1}{2000/02/29}{Added subappendices environment}
%    \begin{macrocode}
\newenvironment{subappendices}{%
  \@resets@ppsub
%    \end{macrocode}
% There are two options applicable to the chapter style. To implement
% the \Lopt{titletoc} option, we redefine the |\addcontentsline| command.
%    \begin{macrocode}
  \if@chapter@pp
    \if@dotitletoc@pp \@redotocentry@pp{section} \fi
%    \end{macrocode}
% To implement the \Lopt{title} option we do cunning things with the
% |\@seccntformat| command.
%    \begin{macrocode}
    \if@dotitle@pp
      \def\sectionname{\appendixname}
      \def\@seccntformat##1{\@ifundefined{##1name}{}{\csname ##1name\endcsname\ }%
        \csname the##1\endcsname\quad}
    \fi
  \else
%    \end{macrocode}
% The rest of the code is for the section style.
%    \begin{macrocode}
    \if@dotitletoc@pp \@redotocentry@pp{subsection} \fi
    \if@dotitle@pp
      \def\subsectionname{\appendixname}
      \def\@seccntformat##1{\@ifundefined{##1name}{}{\csname ##1name\endcsname\ }%
        \csname the##1\endcsname\quad}
    \fi
  \fi}{}

%    \end{macrocode}
% \end{environment}
%
% \begin{macro}{\@formatsecmark@pp}
% Formats the page header for a redefined |\sectionmark|.
%    \begin{macrocode}
\newcommand{\@formatsecmark@pp}[1]{%
  \MakeUppercase{\appendixname\space
    \ifnum \c@secnumdepth >\z@
      \thesection\quad
    \fi
    #1}}
%    \end{macrocode}
% \end{macro}
%
% \begin{macro}{\@redotocentry@pp}
% In order to implement the \Lopt{titletoc} option we redefine the
% |\addcontentsline| command which is used to put entries into the ToC.
% |\@redotocentry@pp{|\meta{sect}|}| does the redefinition, where \meta{sect}
% is the name of the sectional heading (i.e., either chapter or section).
% \changes{v1.2}{2002/08/06}{Replaced ifthen package code in 
%                \cs{@redotocentry@pp} by \cs{ifx} code.}
% \changes{v1.2}{2002/08/06}{HW mods to \cs{@redotocentry@pp}}
%    \begin{macrocode}
\newcommand{\@redotocentry@pp}[1]{%
%    \end{macrocode}
% Save the original definition of |\addcontentsline|. Then start the 
% redefinition.
%    \begin{macrocode}
  \let\oldacl@pp=\addcontentsline
  \def\addcontentsline##1##2##3{%
%    \end{macrocode}
% Check if writing to ToC and appropriate section.
%    \begin{macrocode}
    \def\@pptempa{##1}\def\@pptempb{toc}%
    \ifx\@pptempa\@pptempb
%    \end{macrocode}
% Adding to the ToC file, so check on the sectioning command.
%    \begin{macrocode}
      \def\@pptempa{##2}\def\@pptempb{#1}%
      \ifx\@pptempa\@pptempb
%    \end{macrocode}
% The sectioning command is the same as that specified by the argument to 
% |\@redotocentry@pp|, so get on with the redefinition.
%    \begin{macrocode}
	\oldacl@pp{##1}{##2}{\appendixname\space ##3}%
      \else
%    \end{macrocode}
% The heading was different from the argument. No redefinition is required, so
% call the original |\addcontentsline|.
%    \begin{macrocode}
        \oldacl@pp{##1}{##2}{##3}%
      \fi
    \else
%    \end{macrocode}
% Adding to a file that is not the ToC. No redefinition is required, so
% call the original |\addcontentsline|.
%    \begin{macrocode}
      \oldacl@pp{##1}{##2}{##3}%
    \fi}
}
%    \end{macrocode}
% \end{macro}
%
%
%
%    The end of this package.
%    \begin{macrocode}
%</usc>
%    \end{macrocode}
%
%
% \bibliographystyle{alpha}
%
% \begin{thebibliography}{GMS94}
%
% \bibitem[GMS94]{GOOSSENS94}
% Michel Goossens, Frank Mittelbach, and Alexander Samarin.
% \newblock {\em The LaTeX Companion}.
% \newblock Addison-Wesley Publishing Company, 1994.
%
% \bibitem[Wil96]{PRW96i}
% Peter~R. Wilson.
% \newblock {\em {LaTeX for standards: The LaTeX package files user manual}}.
% \newblock NIST Report NISTIR, June 1996.
%
% \end{thebibliography}
%
% \PrintIndex
% \Finale
%
\endinput

%% \CharacterTable
%%  {Upper-case    \A\B\C\D\E\F\G\H\I\J\K\L\M\N\O\P\Q\R\S\T\U\V\W\X\Y\Z
%%   Lower-case    \a\b\c\d\e\f\g\h\i\j\k\l\m\n\o\p\q\r\s\t\u\v\w\x\y\z
%%   Digits        \0\1\2\3\4\5\6\7\8\9
%%   Exclamation   \!     Double quote  \"     Hash (number) \#
%%   Dollar        \$     Percent       \%     Ampersand     \&
%%   Acute accent  \'     Left paren    \(     Right paren   \)
%%   Asterisk      \*     Plus          \+     Comma         \,
%%   Minus         \-     Point         \.     Solidus       \/
%%   Colon         \:     Semicolon     \;     Less than     \<
%%   Equals        \=     Greater than  \>     Question mark \?
%%   Commercial at \@     Left bracket  \[     Backslash     \\
%%   Right bracket \]     Circumflex    \^     Underscore    \_
%%   Grave accent  \`     Left brace    \{     Vertical bar  \|
%%   Right brace   \}     Tilde         \~}


